\textbf{Operaciones para comparar ambos casos.}
\begin{enumerate}

\item Contribución cinética del caudal.

\[\fitmath{\mathrm{Fk}=2.4\cdot 2500\approx 6000\,\mathrm{N}}\]

\item Presión residual en el primer ambiente.

\[\fitmath{\mathrm{Fp}_{90}=\left(45000-90000\right)\cdot 0.012\approx -540\,\mathrm{N}}\]

\item Sumar contribuciones del mismo balance.

\[\fitmath{F_{90}=\mathrm{Fk}+\mathrm{Fp}_{90}}\]

\[\fitmath{F_{90}=(6000)+(-540)\approx 5460\,\mathrm{N}}\]

\item Velocidad efectiva incluye presión.

\[\fitmath{c_{90}=\frac{F_{90}}{2.4}}\]

\[\fitmath{c_{90}=\frac{(5460)}{2.4}\approx 2275\,\mathrm{m/s}}\]

\item Pasar a impulso específico convencional.

\[\fitmath{\mathrm{Isp}_{90}=\frac{c_{90}}{9.80665}}\]

\[\fitmath{\mathrm{Isp}_{90}=\frac{(2275)}{9.80665}\approx 231.9854\,\mathrm{s}}\]

\item En vacío se resta presión ambiente cero.

\[\fitmath{\mathrm{Fv}=\mathrm{Fk}+45000\cdot 0.012}\]

\[\fitmath{\mathrm{Fv}=(6000)+45000\cdot 0.012\approx 6540\,\mathrm{N}}\]

\item Velocidad efectiva de vacío.

\[\fitmath{\mathrm{cv}=\frac{\mathrm{Fv}}{2.4}}\]

\[\fitmath{\mathrm{cv}=\frac{(6540)}{2.4}\approx 2725\,\mathrm{m/s}}\]

\item Impulso específico de vacío.

\[\fitmath{\mathrm{Ispv}=\frac{\mathrm{cv}}{9.80665}}\]

\[\fitmath{\mathrm{Ispv}=\frac{(2725)}{9.80665}\approx 277.8727\,\mathrm{s}}\]

\end{enumerate}
